\documentclass[10pt,twocolumn]{article}

\usepackage[margin=1.7cm,columnsep=0.65cm]{geometry}
\usepackage[T1]{fontenc}
\usepackage[utf8]{inputenc}
\usepackage[spanish,es-nodecimaldot]{babel}
\usepackage{lmodern}
\usepackage{microtype}
\usepackage{amsmath,amssymb,bm}
\usepackage{graphicx}
\usepackage{booktabs,tabularx,array,multirow}
\usepackage{siunitx}
\usepackage{algorithm}
\usepackage{algpseudocode}
\usepackage[font=small,labelfont=bf]{caption}
\usepackage{subcaption}
\usepackage{xcolor}
\usepackage[hidelinks]{hyperref}

% Dos columnas sin regla de separación y texto completamente negro.
\setlength{\columnseprule}{0pt}
\color{black}
\hypersetup{colorlinks=false}
\graphicspath{{figures/}}
\sisetup{output-decimal-marker={.},round-mode=places,round-precision=4}
\captionsetup[table]{position=top,skip=4pt}
\captionsetup[figure]{position=bottom,skip=4pt}
% Babel usa «Cuadro» por defecto; se redefine dentro del idioma español
% para que el cambio permanezca activo después de \begin{document}.
\addto\captionsspanish{\renewcommand{\tablename}{Tabla}}
\floatname{algorithm}{Algoritmo}
\renewcommand{\algorithmicrequire}{\textbf{Entrada:}}
\renewcommand{\algorithmicensure}{\textbf{Salida:}}
\setlength{\parindent}{1em}
\setlength{\parskip}{0pt}
\raggedbottom
\renewcommand{\textfraction}{0.05}

\begin{document}
\twocolumn[
\begin{center}
{\Large\bfseries
Práctica 2: Clasificación por Funciones de Distancia --- K-means e\\
ISODATA\par}
\vspace{0.7em}
\begin{tabular*}{\textwidth}{@{}l@{\extracolsep{\fill}}r@{}}
Rafael Ortega de la Paz --- No. de cuenta 420054085 &
Sebastian Andre Fuentes Poma --- No. de cuenta 620271005
\end{tabular*}
\vspace{0.25em}

{\itshape Facultad de Ingeniería, Universidad Nacional Autónoma de México
--- Reconocimiento de Patrones\par}
\end{center}
\vspace{0.45em}

\noindent{\small\bfseries Resumen---Este reporte documenta la implementación
y comparación de dos algoritmos de clasificación no supervisada basados en
distancias: K-means e ISODATA. Ambos se implementaron desde cero con NumPy,
incluyendo inicialización aleatoria y K-means++ para K-means, y eliminación,
división y fusión de grupos para ISODATA. Se estudió el efecto de la semilla
sobre la convergencia en datos sintéticos y se compararon ambos métodos mediante
Silhouette, el índice de Davies--Bouldin (DBI) y el índice de Rand ajustado
(ARI). También se analizaron los datos Wine tras estandarizar sus 13
características y aplicar PCA. En Wine, ISODATA identificó tres grupos y obtuvo
resultados ligeramente superiores a K-means en Silhouette y ARI, aunque no en
DBI.\par}
\vspace{1.1em}
]

\section{Introducción y objetivo}

El objetivo de esta práctica es implementar y analizar K-means e ISODATA,
comprender el efecto de sus parámetros y comparar su desempeño en conjuntos
de datos sintéticos y reales. Además de visualizar la evolución de los
centroides y la asignación de grupos, se emplean métricas internas y externas
para evaluar cuantitativamente las particiones obtenidas.

\section{Metodología}

Sea $X=\{\bm{x}_1,\ldots,\bm{x}_n\}\subset\mathbb{R}^p$. Ambos métodos usan
la distancia euclidiana
\begin{equation}
 d(\bm{x}_i,\bm{\mu}_j)
 =\lVert\bm{x}_i-\bm{\mu}_j\rVert_2
 =\sqrt{\sum_{\ell=1}^{p}(x_{i\ell}-\mu_{j\ell})^2}.
 \label{eq:distancia}
\end{equation}

\subsection{K-means}

K-means divide las observaciones en $k$ grupos minimizando la suma de
cuadrados intra-grupo o inercia:
\begin{equation}
 J=\sum_{j=1}^{k}\sum_{\bm{x}_i\in C_j}
 \lVert\bm{x}_i-\bm{\mu}_j\rVert_2^2.
 \label{eq:inercia}
\end{equation}
En la etapa de asignación,
$c_i=\arg\min_j d(\bm{x}_i,\bm{\mu}_j)$; después, cada centroide se actualiza
como
\begin{equation}
 \bm{\mu}_j=\frac{1}{|C_j|}\sum_{\bm{x}_i\in C_j}\bm{x}_i.
 \label{eq:centroide}
\end{equation}
La inicialización aleatoria toma $k$ observaciones del conjunto. En K-means++,
una vez elegidos algunos centroides, el siguiente punto se selecciona con
probabilidad
\begin{equation}
 \Pr(\bm{x}_i)=\frac{D(\bm{x}_i)^2}
 {\sum_{r=1}^{n}D(\bm{x}_r)^2},
 \quad D(\bm{x})=\min_j d(\bm{x},\bm{\mu}_j).
 \label{eq:kpp}
\end{equation}
Se detiene cuando
$\max_j\lVert\bm{\mu}_j^{(t)}-\bm{\mu}_j^{(t-1)}\rVert_2<\varepsilon$.

\begin{algorithm}[H]
\small
\caption{K-means}
\label{alg:kmeans}
\begin{algorithmic}[1]
\Require $X$, $k$, método de inicio, $T$, $\varepsilon$
\Ensure Etiquetas $\bm{c}$ y centroides $\{\bm{\mu}_j\}$
\State Inicializar $k$ centroides (aleatorio o K-means++)
\For{$t=1,\ldots,T$}
  \State $c_i\gets\arg\min_j\lVert\bm{x}_i-\bm{\mu}_j\rVert_2$
  \State $\bm{\mu}'_j\gets |C_j|^{-1}\sum_{\bm{x}_i\in C_j}\bm{x}_i$
  \If{$\max_j\lVert\bm{\mu}'_j-\bm{\mu}_j\rVert_2<\varepsilon$}
    \State \textbf{terminar}
  \EndIf
  \State $\bm{\mu}_j\gets\bm{\mu}'_j$
\EndFor
\end{algorithmic}
\end{algorithm}

\subsection{ISODATA}

ISODATA extiende K-means y modifica dinámicamente el número de grupos. Para
cada $C_j$ se calculan su tamaño $n_j$, su centroide y su desviación estándar
por componente:
\begin{equation}
 \sigma_{j\ell}=\sqrt{\frac{1}{n_j}
 \sum_{\bm{x}_i\in C_j}(x_{i\ell}-\mu_{j\ell})^2}.
 \label{eq:sigma}
\end{equation}
Los grupos con $n_j<N_{\min}$ se eliminan. Si
$\max_\ell\sigma_{j\ell}>\sigma_{\max}$ y $k<k_{\max}$, se divide el grupo
sobre la dirección $\ell^*$ de mayor dispersión:
\begin{equation}
 \bm{\mu}_{j}^{\pm}=\bm{\mu}_j
 \pm \alpha\sigma_{j\ell^*}\bm{e}_{\ell^*}.
 \label{eq:division}
\end{equation}
Si $d(\bm{\mu}_a,\bm{\mu}_b)<d_{\min}$, los grupos se fusionan mediante el
promedio ponderado
\begin{equation}
 \bm{\mu}_{ab}=\frac{n_a\bm{\mu}_a+n_b\bm{\mu}_b}{n_a+n_b}.
 \label{eq:fusion}
\end{equation}

\begin{algorithm}[H]
\small
\caption{ISODATA}
\label{alg:isodata}
\begin{algorithmic}[1]
\Require $X$, $k_{\mathrm{init}}$, $N_{\min}$, $\sigma_{\max}$,
$d_{\min}$, $\theta_c$, $k_{\max}$, $T$, $\varepsilon$
\Ensure Etiquetas y centroides finales
\State Inicializar $k_{\mathrm{init}}$ centroides
\For{$t=1,\ldots,T$}
  \State Asignar cada punto al centroide más cercano
  \State Eliminar $C_j$ si $n_j<N_{\min}$ y reasignar sus puntos
  \State Recalcular centroides y desviaciones $\sigma_{j\ell}$
  \If{existe $\max_\ell\sigma_{j\ell}>\sigma_{\max}$ y $k<k_{\max}$}
    \State Dividir $C_j$ con la Ec.~\eqref{eq:division}
  \Else
    \State Fusionar hasta $\theta_c$ pares con distancia $<d_{\min}$
  \EndIf
  \If{el desplazamiento máximo es $<\varepsilon$}
    \State \textbf{terminar}
  \EndIf
\EndFor
\end{algorithmic}
\end{algorithm}

\subsection{Preprocesamiento y métricas}

Para Wine, cada variable se estandarizó como
$z_{i\ell}=(x_{i\ell}-\bar{x}_{\ell})/s_{\ell}$. PCA diagonaliza la matriz de
covarianza $S=V\Lambda V^\top$ y proyecta los datos mediante $Z_m=ZV_m$.
Se eligió el menor $m$ que satisface
\begin{equation}
 \frac{\sum_{r=1}^{m}\lambda_r}{\sum_{r=1}^{p}\lambda_r}\geq0.95.
 \label{eq:pca}
\end{equation}
Entre las tres estrategias de preprocesamiento consideradas para Wine
--usar las 13 características estandarizadas directamente, reducir
mediante PCA, o seleccionar variables por correlación de Pearson-- se
optó por PCA porque obtuvo, en pruebas preliminares del equipo, mejor
Silhouette y menor DBI que usar las 13 características sin reducir, y
porque la selección por Pearson con un umbral razonable no eliminó
variables redundantes de forma efectiva en este conjunto de datos. PCA
además mitiga el efecto de la maldición de la dimensionalidad sobre las
distancias euclidianas que emplean ambos algoritmos.

Para cada observación, el coeficiente Silhouette es
\begin{equation}
 s_i=\frac{b_i-a_i}{\max(a_i,b_i)},
 \label{eq:silhouette}
\end{equation}
donde $a_i$ es la distancia media al propio grupo y $b_i$ la menor distancia
media a otro grupo. El DBI se define como
\begin{equation}
 \mathrm{DBI}=\frac{1}{k}\sum_{i=1}^{k}\max_{j\ne i}
 \frac{S_i+S_j}{d(\bm{\mu}_i,\bm{\mu}_j)},
 \label{eq:dbi}
\end{equation}
donde $S_i$ mide la dispersión de $C_i$. Silhouette y ARI se prefieren altos;
DBI se prefiere bajo. El ARI corrige por azar la concordancia entre las
etiquetas estimadas y las reales.

\section{Resultados}

\subsection{K-means en datos sintéticos}

Primero se generaron 400 observaciones bidimensionales distribuidas en cuatro
grupos mediante \texttt{make\_blobs}, con \texttt{cluster\_std=1.7} y
\texttt{random\_state=20}. La Figura~\ref{fig:sinteticos} muestra los datos
antes de aplicar cualquier algoritmo; los colores corresponden a las clases
reales y sólo se usan como referencia para evaluar el agrupamiento.

\begin{figure}[H]
  \centering
  \includegraphics[width=\columnwidth]{datos_sinteticos.png}
  \caption{Datos sintéticos de entrada: 400 observaciones, dos características
  y cuatro clases reales.}
  \label{fig:sinteticos}
\end{figure}

K-means se ejecutó con $k=4$ usando las semillas 10, 20 y 30 para cada método
de inicialización. La Figura~\ref{fig:kmeans-corridas} reúne las seis
particiones finales: la fila superior corresponde al inicio aleatorio y la
inferior a K-means++. Cada color representa un grupo estimado y las cruces
marcan los centroides finales. Esta comparación permite observar la
sensibilidad del algoritmo a la posición inicial de los centroides.

\begin{figure}[H]
  \centering
  \includegraphics[width=\columnwidth]{kmeans_corridas.png}
  \caption{K-means: asignaciones y centroides finales para tres semillas con
  inicio aleatorio (arriba) y K-means++ (abajo).}
  \label{fig:kmeans-corridas}
\end{figure}

La Figura~\ref{fig:kmeans-inercia} presenta la inercia después de cada
iteración. Una curva descendente indica que las actualizaciones mejoran la
partición; la meseta señala la convergencia. Cinco corridas alcanzan el mismo
mínimo, mientras que K-means++ con semilla 20 se estabiliza en un valor mayor.

\begin{figure}[H]
  \centering
  \includegraphics[width=\columnwidth]{kmeans_inercia.png}
  \caption{K-means: evolución de la inercia en cada una de las seis corridas.}
  \label{fig:kmeans-inercia}
\end{figure}

% La tabla y su interpretación se mantienen juntas y en ese orden.
\noindent\begin{minipage}{\columnwidth}
{\centering
\captionsetup{hypcap=false}
\captionof{table}{Inercia final de K-means ($k=4$).}
\label{tab:inercia-kmeans}
\small
\setlength{\tabcolsep}{3.5pt}
\begin{tabular}{lrrrr}
\toprule
Inicio & Mín. & Máx. & Media & Rango\\
\midrule
Aleatorio & 2086.69 & 2086.69 & 2086.69 & 0.00\\
K-means++ & 2086.69 & 3679.84 & 2617.79 & 1593.15\\
\bottomrule
\end{tabular}
\par}
\smallskip

La Tabla~\ref{tab:inercia-kmeans} cuantifica lo observado en las figuras. Las
tres corridas aleatorias alcanzaron el mejor valor, mientras que K-means++
quedó en un mínimo local con la semilla 20 (inercia 76\% mayor). K-means++
reduce probabilísticamente el riesgo de una mala inicialización, pero no lo
elimina; con sólo tres repeticiones, una corrida atípica modifica fuertemente
el promedio.
\end{minipage}

\subsection{ISODATA en datos sintéticos}

ISODATA se ejecutó con $k_{\mathrm{init}}=6$ y $k_{\mathrm{init}}=3$. El primer
caso convergió en seis iteraciones a cuatro grupos e inercia 2086.6895; el
segundo lo hizo en cinco iteraciones, también con cuatro grupos, e inercia
2086.8061. Así, el método recuperó el número real de grupos desde una
sobreestimación y desde una subestimación inicial. La
Figura~\ref{fig:isodata-k} registra el número de grupos en cada iteración: con
$k_{\mathrm{init}}=6$ se eliminan o fusionan grupos hasta llegar a cuatro;
con $k_{\mathrm{init}}=3$ se divide un grupo disperso para alcanzar el mismo
valor.

\begin{figure}[H]
  \centering
  \includegraphics[width=\columnwidth]{isodata_k_iteracion.png}
  \caption{ISODATA: evolución del número de grupos desde
  $k_{\mathrm{init}}=6$ y 3; la línea punteada marca el valor real $k=4$.}
  \label{fig:isodata-k}
\end{figure}

La Figura~\ref{fig:isodata-inercia} verifica que ambas configuraciones reducen
rápidamente el error intra-grupo y terminan en valores casi idénticos.

\begin{figure}[H]
  \centering
  \includegraphics[width=\columnwidth]{isodata_inercia.png}
  \caption{ISODATA: inercia por iteración para los dos valores de
  $k_{\mathrm{init}}$.}
  \label{fig:isodata-inercia}
\end{figure}

La Figura~\ref{fig:isodata-clusters} muestra el resultado espacial final. Los
dos puntos de partida producen cuatro regiones y centroides prácticamente
coincidentes, lo que confirma que la adaptación de $k$ no alteró la solución
óptima de forma apreciable.

\begin{figure}[H]
  \centering
  \includegraphics[width=\columnwidth]{isodata_clusters.png}
  \caption{ISODATA: grupos y centroides finales desde
  $k_{\mathrm{init}}=6$ (izquierda) y 3 (derecha).}
  \label{fig:isodata-clusters}
\end{figure}

La comparación numérica con la mejor corrida de K-means se presenta en la
Tabla~\ref{tab:comparacion-sintetica}.

\begin{table}[H]
\centering
\caption{Comparación de las soluciones sintéticas.}
\label{tab:comparacion-sintetica}
\small
\begin{tabular}{lcc}
\toprule
Algoritmo & Grupos & Inercia\\
\midrule
K-means (mejor) & 4 & 2086.6895\\
ISODATA ($k_{\mathrm{init}}=6$) & 4 & 2086.6895\\
ISODATA ($k_{\mathrm{init}}=3$) & 4 & 2086.8061\\
\bottomrule
\end{tabular}
\end{table}

Las soluciones coinciden en posición con diferencias menores que 0.04 por
coordenada. ISODATA, por tanto, alcanza prácticamente la misma partición que
K-means sin requerir que $k$ sea conocido de antemano.

\subsection{Métricas en datos sintéticos}

Para el Ejercicio 3 se eligió la corrida de menor inercia de cada algoritmo.
La Tabla~\ref{tab:metricas-sinteticas} compara las tres métricas solicitadas.

\begin{table}[H]
\centering
\caption{Métricas en los datos sintéticos.}
\label{tab:metricas-sinteticas}
\small
\begin{tabular}{lccc}
\toprule
Algoritmo & Silhouette $\uparrow$ & DBI $\downarrow$ & ARI $\uparrow$\\
\midrule
K-means & 0.5639 & 0.6131 & 0.9156\\
ISODATA & 0.5639 & 0.6131 & 0.9156\\
\bottomrule
\end{tabular}
\end{table}

Las métricas de la Tabla~\ref{tab:metricas-sinteticas} son idénticas porque
las mejores corridas produjeron la misma asignación. El Silhouette de 0.56 y
el DBI de 0.61 indican grupos compactos y razonablemente separados; el ARI de
0.92 muestra alta concordancia con las etiquetas reales.

\subsection{Conjunto de datos Wine}

Wine contiene 178 observaciones, 13 características con escalas distintas y
tres clases conocidas. Tras estandarizar, PCA retuvo 10 componentes, que
explican 96.17\% de la varianza. La Figura~\ref{fig:wine-pca} muestra la
varianza acumulada y justifica la selección: el umbral de 95\% se supera por
primera vez al conservar diez componentes.

La Tabla~\ref{tab:eigenvalores-wine} presenta los eigenvalores de la matriz de
covarianza de los datos estandarizados. Los primeros componentes concentran la
mayor parte de la información: PC1 y PC2 explican conjuntamente 55.41\% de la
varianza, mientras que los diez componentes retenidos acumulan 96.17\%.

\begin{table}[H]
\centering
\caption{Eigenvalores y varianza explicada del PCA de Wine.}
\label{tab:eigenvalores-wine}
\small
\setlength{\tabcolsep}{4pt}
\begin{tabular}{crrr}
\toprule
PC & Eigenvalor $\lambda_j$ & Var. (\%) & Acum. (\%)\\
\midrule
1  & 4.7324 & 36.20 & 36.20\\
2  & 2.5111 & 19.21 & 55.41\\
3  & 1.4542 & 11.12 & 66.53\\
4  & 0.9242 &  7.07 & 73.60\\
5  & 0.8580 &  6.56 & 80.16\\
6  & 0.6453 &  4.94 & 85.10\\
7  & 0.5541 &  4.24 & 89.34\\
8  & 0.3505 &  2.68 & 92.02\\
9  & 0.2905 &  2.22 & 94.24\\
10 & 0.2523 &  1.93 & 96.17\\
11 & 0.2271 &  1.74 & 97.91\\
12 & 0.1697 &  1.30 & 99.20\\
13 & 0.1040 &  0.80 & 100.00\\
\bottomrule
\end{tabular}
\end{table}

Los coeficientes de los eigenvectores asociados a los diez componentes
conservados aparecen en la Tabla~\ref{tab:eigenvectores-wine}. Cada columna es
un eigenvector o componente principal, mientras que las filas indican sus
coordenadas respecto a las 13 características originales estandarizadas. La
magnitud absoluta de cada coeficiente indica el peso de la característica y su
signo determina la dirección de la asociación.

\begin{table*}[t]
\centering
\caption{Coeficientes de los eigenvectores del PCA de Wine.}
\label{tab:eigenvectores-wine}
\footnotesize
\setlength{\tabcolsep}{4.5pt}
\renewcommand{\arraystretch}{1.05}
\resizebox{\textwidth}{!}{%
\begin{tabular}{lrrrrrrrrrr}
\toprule
Característica & PC1 & PC2 & PC3 & PC4 & PC5 & PC6 & PC7 & PC8 & PC9 & PC10\\
\midrule
Alcohol                        &  0.144 &  0.484 & -0.207 & -0.018 & -0.266 &  0.214 & -0.056 &  0.396 & -0.509 &  0.212\\
Ácido málico                   & -0.245 &  0.225 &  0.089 &  0.537 &  0.035 &  0.537 &  0.421 &  0.066 &  0.075 & -0.309\\
Ceniza                         & -0.002 &  0.316 &  0.626 & -0.214 & -0.143 &  0.154 & -0.149 & -0.170 &  0.308 & -0.027\\
Alcalinidad de la ceniza       & -0.239 & -0.011 &  0.612 &  0.061 &  0.066 & -0.101 & -0.287 &  0.428 & -0.200 &  0.053\\
Magnesio                       &  0.142 &  0.300 &  0.131 & -0.352 &  0.727 &  0.038 &  0.323 & -0.156 & -0.271 &  0.068\\
Fenoles totales                &  0.395 &  0.065 &  0.146 &  0.198 & -0.149 & -0.084 & -0.028 & -0.406 & -0.286 & -0.320\\
Flavonoides                    &  0.423 & -0.003 &  0.151 &  0.152 & -0.109 & -0.019 & -0.061 & -0.187 & -0.050 & -0.163\\
Fenoles no flavonoides         & -0.299 &  0.029 &  0.170 & -0.203 & -0.501 & -0.259 &  0.595 & -0.233 & -0.196 &  0.216\\
Proantocianinas                &  0.313 &  0.039 &  0.149 &  0.399 &  0.137 & -0.534 &  0.372 &  0.368 &  0.209 &  0.134\\
Intensidad de color            & -0.089 &  0.530 & -0.137 &  0.066 & -0.076 & -0.419 & -0.228 & -0.034 & -0.056 & -0.291\\
Tono                           &  0.297 & -0.279 &  0.085 & -0.428 & -0.174 &  0.106 &  0.232 &  0.437 & -0.086 & -0.522\\
OD280/OD315 de vinos diluidos  &  0.376 & -0.164 &  0.166 &  0.184 & -0.101 &  0.266 & -0.045 & -0.078 & -0.137 &  0.524\\
Prolina                        &  0.287 &  0.365 & -0.127 & -0.232 & -0.158 &  0.120 &  0.077 &  0.120 &  0.576 &  0.162\\
\bottomrule
\end{tabular}%
}
\end{table*}

\begin{figure}[H]
  \centering
  \includegraphics[width=\columnwidth]{wine_pca.png}
  \caption{Wine: varianza explicada acumulada por PCA y umbral de 95\%.}
  \label{fig:wine-pca}
\end{figure}

K-means se ejecutó con $k=3$, K-means++ y semilla 20; convergió en nueve
iteraciones con inercia 1189.64. Para ISODATA se usaron
$k_{\mathrm{init}}=5$, $N_{\min}=8$, $\sigma_{\max}=1.6$,
$d_{\min}=2.2$, $\theta_c=2$ y $k_{\max}=6$; estos valores se calibraron
a la escala de las diez componentes principales, cuya varianza es del
orden de 1 tras estandarizar, de modo que $\sigma_{\max}$ y $d_{\min}$
se fijaron en ese mismo orden de magnitud en lugar de usar los valores
de referencia del caso sintético, que estaba en la escala original de
los datos sin estandarizar. ISODATA convergió a tres grupos
con inercia 1190.36. La Figura~\ref{fig:wine-convergencia} coloca las curvas
de inercia junto a la evolución de $k$: las dos inercias se estabilizan en
valores similares e ISODATA reduce el número de grupos de cinco a tres.

\begin{figure}[H]
  \centering
  \includegraphics[width=\columnwidth]{wine_convergencia.png}
  \caption{Wine: inercia por iteración de ambos algoritmos y evolución del
  número de grupos de ISODATA.}
  \label{fig:wine-convergencia}
\end{figure}

La Tabla~\ref{tab:centroides-wine} muestra las coordenadas de los centroides
finales en las diez componentes principales utilizadas para el agrupamiento.
Los resultados aparecen en un orden distinto porque, al ser algoritmos no
supervisados, los números asignados a los grupos son arbitrarios y no indican
una jerarquía. Por ello, el grupo 0 de K-means corresponde al grupo 2 de
ISODATA, aunque ambos representan prácticamente el mismo conjunto.

\begin{table*}[t]
\centering
\caption{Centroides finales de K-means e ISODATA para Wine en el espacio PCA.}
\label{tab:centroides-wine}
\footnotesize
\setlength{\tabcolsep}{4pt}
\renewcommand{\arraystretch}{1.1}
\resizebox{\textwidth}{!}{%
\begin{tabular}{llrrrrrrrrrr}
\toprule
Algoritmo & Grupo & PC1 & PC2 & PC3 & PC4 & PC5 & PC6 & PC7 & PC8 & PC9 & PC10\\
\midrule
\multirow{3}{*}{K-means}
 & 0 & -2.720 &  1.126 & -0.239 &  0.063 &  0.074 & -0.100 & -0.060 &  0.007 & -0.020 & -0.062\\
 & 1 & -0.037 & -1.772 &  0.186 &  0.080 &  0.071 & -0.130 & -0.002 &  0.018 & -0.032 & -0.023\\
 & 2 &  2.276 &  0.932 &  0.002 & -0.136 & -0.135 &  0.218 &  0.052 & -0.025 &  0.050 &  0.075\\
\midrule
\multirow{3}{*}{ISODATA}
 & 0 &  2.292 &  0.943 & -0.086 & -0.142 & -0.141 &  0.235 &  0.063 &  0.004 &  0.058 &  0.080\\
 & 1 & -0.016 & -1.742 &  0.264 &  0.083 &  0.073 & -0.140 & -0.011 & -0.009 & -0.038 & -0.027\\
 & 2 & -2.720 &  1.126 & -0.239 &  0.063 &  0.074 & -0.100 & -0.060 &  0.007 & -0.020 & -0.062\\
\bottomrule
\end{tabular}%
}
\end{table*}

La Tabla~\ref{tab:contingencia-wine} cuenta cuántas observaciones de cada clase
real contiene cada grupo. Una fila concentrada en una sola columna indica una
buena correspondencia entre grupo y clase.

\begin{table}[H]
\centering
\caption{Matrices de contingencia para Wine.}
\label{tab:contingencia-wine}
\begin{subtable}[t]{0.98\columnwidth}
\centering
\caption{K-means ($k=3$)}
\begin{tabular}{lrrr}
\toprule
Grupo & Clase 1 & Clase 2 & Clase 3\\
\midrule
0 & 0 & 3 & 48\\
1 & 0 & 65 & 0\\
2 & 59 & 3 & 0\\
\bottomrule
\end{tabular}
\end{subtable}
\medskip

\begin{subtable}[t]{0.98\columnwidth}
\centering
\caption{ISODATA ($k=3$)}
\begin{tabular}{lrrr}
\toprule
Grupo & Clase 1 & Clase 2 & Clase 3\\
\midrule
0 & 59 & 2 & 0\\
1 & 0 & 66 & 0\\
2 & 0 & 3 & 48\\
\bottomrule
\end{tabular}
\end{subtable}
\end{table}

Cada grupo está dominado casi completamente por una clase real. El principal
traslape ocurre entre las clases 2 y 3. Como apoyo visual, la
Figura~\ref{fig:wine-clusters} proyecta la solución de K-means sobre PC1 y PC2:
las clases asociadas a los grupos morado y amarillo se separan claramente,
mientras que el grupo turquesa presenta mayor dispersión y cercanía con ambos.

\begin{figure}[H]
  \centering
  \includegraphics[width=\columnwidth]{wine_clusters.png}
  \caption{Wine: proyección PC1--PC2 coloreada según el grupo asignado por
  K-means. Esta figura es una visualización; el agrupamiento se calculó con
  las 10 componentes retenidas.}
  \label{fig:wine-clusters}
\end{figure}

Finalmente, la Tabla~\ref{tab:metricas-wine} resume la calidad global de ambas
particiones en las diez componentes.

\begin{table}[H]
\centering
\caption{Métricas de evaluación en Wine.}
\label{tab:metricas-wine}
\small
\begin{tabular}{lccc}
\toprule
Algoritmo & Silhouette $\uparrow$ & DBI $\downarrow$ & ARI $\uparrow$\\
\midrule
K-means & 0.2987 & \textbf{1.3363} & 0.8975\\
ISODATA & \textbf{0.2997} & 1.3387 & \textbf{0.9149}\\
\bottomrule
\end{tabular}
\end{table}

ISODATA obtuvo una mejora pequeña en Silhouette y una mejora más clara en ARI;
K-means presentó un DBI ligeramente menor. La diferencia de DBI es mínima,
por lo que el resultado global es comparable, con ventaja de ISODATA al
recuperar $k=3$ automáticamente y aproximar mejor las clases reales. Para
mejorar ambos métodos podría realizarse una búsqueda sistemática de
$\sigma_{\max}$ y $d_{\min}$, variar el número de componentes o estudiar una
distancia de Mahalanobis, adecuada para grupos no esféricos.

\section{Conclusiones}

La inicialización influye en K-means, pero su ventaja es estadística y no está
garantizada en cada corrida. Con las semillas estudiadas, el inicio aleatorio
superó por azar a K-means++. ISODATA estimó correctamente cuatro grupos en los
datos sintéticos desde $k_{\mathrm{init}}=3$ y 6, y tres grupos en Wine desde
$k_{\mathrm{init}}=5$. En Wine obtuvo el mejor ARI y un Silhouette apenas
superior, mientras K-means logró un DBI apenas menor. La estandarización y PCA
fueron esenciales para trabajar con características de escalas diferentes; en
ISODATA, la calibración de los umbrales a esa nueva escala resultó igualmente
determinante.

\begin{thebibliography}{9}
\small
\bibitem{macqueen}
J. MacQueen, ``Some methods for classification and analysis of multivariate
observations,'' en \emph{Proc. 5th Berkeley Symp. Math. Statist. Prob.}, vol.~1,
1967, pp.~281--297.

\bibitem{kmeanspp}
D. Arthur y S. Vassilvitskii, ``k-means++: The advantages of careful
seeding,'' en \emph{Proc. 18th ACM--SIAM Symp. Discrete Algorithms}, 2007,
pp.~1027--1035.

\bibitem{isodata}
G. H. Ball y D. J. Hall, ``ISODATA, a novel method of data analysis and
pattern classification,'' Stanford Research Institute, Tech. Rep. AD699616,
1965.

\bibitem{sklearn}
F. Pedregosa et al., ``Scikit-learn: Machine learning in Python,''
\emph{J. Mach. Learn. Res.}, vol.~12, pp.~2825--2830, 2011.
\end{thebibliography}

\end{document}
